\chapter{Conclusion}

While distributed systems and operating systems face different requirements, we were able to identify parallisms between the two fields. \\
The most relevant areas of transfer were decomposition, scalable load management and both internal and external diagnosis.

These aspects become especially relevant when taking the collaborative nature of the Linux project into account. Decomposition is already required, in order to be able to assign maintainers to subsystems and enabling specialization. Being resilient towards high loads, both in absolute and relative terms, attracts larger computing firms, and larger financial means. The diagnostic tools (watchdogs and sanitizers) shift safety from trusting potentially inexperienced developers towards reproducable run-time testing.

We conclude that the Linux kernel's resilience strategies are intertwined with its collaborative nature. Its empiric sucess in high stake computing applications should inspire other critical fields - such as energy or infrastructure - to consider open source solutions as a mechanism of standardization and cooperative safety.